\documentclass[10pt,a4paper]{amsart}
\usepackage[margin=19mm]{geometry}
\usepackage{amsmath,amssymb,mathtools}
\usepackage[scr=rsfs,cal=euler]{mathalfa}
\usepackage{xfrac}
\usepackage[unicode,hidelinks,bookmarks=false]{hyperref}
\newcommand{\ie}{i.e.\ }
\newcommand{\cf}{cf.\ }
\newcommand{\resp}{respectively\ }
\newcommand{\Subsection}{\subsection}
\providecommand{\nobreakdash}{\nobreak\hbox{-}}
\setlength{\emergencystretch}{2em}
\multlinegap0pt
\usepackage{amssymb}
\usepackage{mathtools}
\newtheorem{lemma}{Lemma}
\newcommand{\R}{\mathbb R}
\newcommand{\cL}{\mathcal L}
\newcommand{\norm}[2]{\lVert #1\rVert_{#2}}
\title{Strong and Weak-Type Dispersive Estimates for the Energy-Critical Nonlinear Schrödinger Equation with an Inverse-Square Potential}
\author{Tiesong Jiang, Kexue Li}
\date{Source: arXiv:2609.07585v1 (CC BY 4.0)}
\begin{document}
\maketitle

\noindent\emph{Source and licence.} Strong and Weak-Type Dispersive Estimates for the Energy-Critical Nonlinear Schrödinger Equation with an Inverse-Square Potential. Version arXiv:2609.07585v1 (CC BY 4.0), distributed by arXiv under the Creative Commons Attribution 4.0 International licence, http://creativecommons.org/licenses/by/4.0/, as recorded in the arXiv licence metadata for this exact version and shown on the abstract page https://arxiv.org/abs/2609.07585v1. Full licence notice: \texttt{licenses/arxiv-2609.07585-CC-BY-4.0.txt}.\par\smallskip

\noindent\textit{Editorial scope.} Counted unit: the article's lemma lem:endpoint-early-high (source0.tex lines 1159--1175). The article's complete original proof of it is delivered with it (source0.tex lines 1177--1208, one \texttt{proof} environment of 32 lines). No mathematics is written, repaired or re-derived: the statement and the proof are the article's own lines, in the article's own notation, and no retained line is substituted or re-flowed.  Retained uncounted as the source-local context the proof needs: lines 298--303; lines 511--545; lines 547--579; lines 687--693; lines 1087--1102.  Nothing was omitted from any retained span, so the package carries no omission record.  No retained line contains a citation, so the wrapper carries no bibliography and no bibliography entry.  No retained line contains a figure environment, an image-inclusion command or drawing code, so no image is delivered and the package declares no asset dependency.  Editorial formatting only, in addition: an AMS \texttt{amsart} preamble that restates the article's own declarations and macros used by the retained lines, namely the environments \texttt{lemma} on separate counters, together with the standard packages those lines need.  The retained environments print in this document as Lemma 1 in order of appearance, whereas the article prints the same items with the numbering of its own full document.  The complete native source of the article as distributed in the arXiv e-print for this exact version is bundled unchanged as \texttt{sources/209642/source0.tex}.
 \begin{equation}
  \norm{U_a(t)f}{L^{p_*,\infty}}
  \lesssim_{a,p_*}|t|^{-\beta_{p_*}}
  \norm{f}{L^{p_*',1}}.
  \label{eq:linear-weak-endpoint}
 \end{equation}

\begin{lemma}[Adapted Sobolev estimates]
\label{lem:adapted-sobolev}
Let $a>-\frac14$, set
$\sigma=\frac12-\sqrt{\frac14+a}$, and let $0<s<2$. For
$f\in C_c^\infty(\R^3\setminus\{0\})$,
\[
 \norm{|\nabla|^sf}{L^r}
 \lesssim_{a,s,r}\norm{\cL_a^{s/2}f}{L^r}
\]
provided
\begin{equation}
 \frac{s+\sigma}{3}<\frac1r<
 \min\left\{1,\frac{3-\sigma}{3}\right\},
 \label{eq:sobolev-forward-range}
\end{equation}
whereas
\[
 \norm{\cL_a^{s/2}f}{L^r}
 \lesssim_{a,s,r}\norm{|\nabla|^sf}{L^r}
\]
provided
\[
 \max\left\{\frac{s}{3},\frac{\sigma}{3}\right\}<\frac1r<
 \min\left\{1,\frac{3-\sigma}{3}\right\}.
\]
If $1<r<q<\infty$, $1/q=1/r-s/3$, and
\eqref{eq:sobolev-forward-range} holds, then
\begin{equation}
 \norm{f}{L^{q,\theta}}
 \lesssim_{a,s,r,q,\theta}
 \norm{\cL_a^{s/2}f}{L^{r,\theta}},
 \qquad 1\leq\theta\leq\infty.
 \label{eq:adapted-embedding}
\end{equation}
\end{lemma}

\begin{lemma}[Global spacetime estimates]
\label{lem:global-spacetime}
Let $E=\norm{u_0}{\dot H^1}$.
\begin{enumerate}
 \item If $a\geq0$, then
 \begin{equation}
 \begin{split}
  &\norm{u}{L_t^4L_x^\infty}
  +\norm{u}{L_t^{20,10/3}L_x^{15/2}}\\
  &\quad
  +\norm{\cL_a^{1/2}u}{L_t^{8,2}L_x^{12/5,4}}
  +\norm{\cL_a^{1/2}u}{L_t^{8,4}L_x^{12/5,4}}
  \leq C_a(E).
 \end{split}
 \label{eq:nonnegative-spacetime}
 \end{equation}
 If also $u_0\in L^{3/2}$, then
 \[
  \norm{u(t)}{L^3}
  \leq C_a(E)|t|^{-1/2}\norm{u_0}{L^{3/2}},
  \qquad t\neq0.
 \]
 \item Suppose $-\frac14+\frac1{25}<a<0$. If $(q,r)$ is
 nonendpoint admissible, $q>2$, $2<r<6$, $2<\vartheta\leq q$, and
 $2\leq\varphi\leq\infty$, then
 \begin{equation}
  \norm{\cL_a^{1/2}u}{L_t^{q,\vartheta}L_x^{r,\varphi}}
  \leq C_{a,q,r,\vartheta,\varphi}(E).
  \label{eq:attractive-spacetime}
 \end{equation}
 The corresponding strong $L_t^qL_x^r$ estimate also holds.
\end{enumerate}
\end{lemma}

\begin{align}
 \sup_{t\in\R}\norm{S_t(f)-S_t(g)}{L^2}
 &\leq C_{a,E}\norm{f-g}{L^2}, \label{eq:difference-L2}\\
 \norm{S_t(f)-S_t(g)}{L^{7/2}}
 &\leq C_{a,E}|t|^{-9/14}\norm{f-g}{L^{7/5}},
 \qquad t\neq0. \label{eq:difference-seven}
\end{align}

\begin{lemma}[Early-time estimate below the attractive endpoint]
\label{lem:early-high}
Let $6\leq p<p_*$, put
$E=\norm{u_0}{\dot H^1}$ and $A=\norm{u_0}{L^{p'}}$.
Then
\begin{equation}
 \int_0^\infty\norm{|u(s)|^4u(s)}{L^{p'}}\,ds
 \leq C_{a,p}(E)A.
 \label{eq:early-integrability}
\end{equation}
Consequently, for $t>0$,
\[
 \left\|\int_0^{t/2}U_a(t-s)(|u|^4u)(s)\,ds\right\|_{L^p}
 \leq C_{a,p}(E)t^{-\beta_p}A.
\]
\end{lemma}

\begin{lemma}[Early-time estimate at the attractive endpoint]
\label{lem:endpoint-early-high}
Put $E=\norm{u_0}{\dot H^1}$ and
$A_*=\norm{u_0}{L^{p_*',1}}$. Then
\begin{equation}
 \int_0^\infty
 \norm{|u(s)|^4u(s)}{L^{p_*',1}}\,ds
 \leq C_a(E)A_*.
 \label{eq:endpoint-early-integrability}
\end{equation}
Consequently, for $t>0$,
\[
 \left\|\int_0^{t/2}U_a(t-s)(|u|^4u)(s)\,ds
 \right\|_{L^{p_*,\infty}}
 \leq C_a(E)t^{-\beta_*}A_*.
\]
\end{lemma}

\begin{proof}
Let $\alpha_*,k_*,d_*,R_*,r_*,q_*$, and $\eta_*$ be the values at
$p=p_*$ of the quantities in the proof of Lemma~\ref{lem:early-high}.
The same exponent identities hold, and $p_*>10$ gives
\[
 2<r_*<\frac{13}{6}<\frac3{1+\sigma},
 \qquad 2<k_*<q_*,\qquad \eta_*>2.
\]
Thus Lemmas~\ref{lem:adapted-sobolev} and
\ref{lem:global-spacetime} give
\[
 \norm{u}{L_t^{q_*,k_*}L_x^{R_*,\eta_*}}
 \leq C_a(E).
\]
The identity
\[
 \frac57=\frac{1/\alpha_*}{p_*'}
 +\frac{1-1/\alpha_*}{6}
\]
and Lorentz interpolation imply
$\norm{u_0}{L^{7/5}}^{\alpha_*}\leq C_a(E)A_*$.
Combining \eqref{eq:difference-seven}, with the second datum zero,
and spatial Lorentz H\"older yields
\[
 \norm{|u(s)|^4u(s)}{L^{p_*',1}}
 \leq C_a(E)A_*s^{-d_*}
 \norm{u(s)}{L^{R_*,\eta_*}}^{k_*}.
\]
Since $d_*+k_*/q_*=1$, Lorentz H\"older in time proves
\eqref{eq:endpoint-early-integrability}. The Duhamel estimate follows
from \eqref{eq:linear-weak-endpoint} and $|t-s|\simeq t$ on $(0,t/2)$.
\end{proof}

\end{document}
